\SourceLeaf{262}
\Article{27}
\indent If the curved surface is a sphere of radius \(=R\), we have
\[
\begin{gathered}
\alpha=\beta=\gamma=-2f^0=\frac{1}{RR};\quad f''=0,\quad g'=0,\\
6h^0-f^0f^0=0\quad\text{or}\quad k^0=\frac{1}{24R^4}
\end{gathered}
\]
Formula [14] then becomes
\[
A+B+C=\pi+\frac{\sigma}{RR}
\]
This holds exactly, whereas formulas 11—13 give
\[
\begin{aligned}
A^*&=A-\frac{\sigma}{3RR}-\frac{\sigma}{180R^4}(2pp-qq+4qq'-q'q')\\
B^*&=B-\frac{\sigma}{3RR}+\frac{\sigma}{180R^4}(pp-2qq+2qq'+q'q')\\
C^*&=C-\frac{\sigma}{3RR}+\frac{\sigma}{180R^4}(pp+qq+2qq'-2q'q')
\end{aligned}
\]
or, with the same precision,
\[
\begin{aligned}
A^*&=A-\frac{\sigma}{3RR}-\frac{\sigma}{180R^4}(bb+cc-2aa)\\
B^*&=B-\frac{\sigma}{3RR}-\frac{\sigma}{180R^4}(aa+cc-2bb)\\
C^*&=C-\frac{\sigma}{3RR}-\frac{\sigma}{180R^4}(aa+bb-2cc)
\end{aligned}
\]
If quantities of the fourth order are neglected, this yields the familiar theorem first proposed by the distinguished \textsc{Legendre}.

\Article{28}
\indent When the terms of the fourth order are dropped, our general formulas become very simple. They are
\[
\begin{aligned}
A^*&=A-\tfrac{1}{12}\sigma(2\alpha+\beta+\gamma)\\
B^*&=B-\tfrac{1}{12}\sigma(\alpha+2\beta+\gamma)\\
C^*&=C-\tfrac{1}{12}\sigma(\alpha+\beta+2\gamma)
\end{aligned}
\]
Different corrections must therefore be applied to the angles \(A,B,C\) on a non-spherical surface so that the sines of the resulting angles are proportional to their opposite sides. In general, the differences between the corrections will be of the third order. If the surface differs little from a sphere, these differences will be of a higher order. Even in the largest triangles on the Earth’s surface whose angles can actually be measured, the difference can always be regarded as
